% GENERATED FILE. Edit canonical lessons, then run npm run generate:books.
% canonical-source-hash: fnv1a64:cdcd5e9b9f019aec
% canonical-lessons: GU-C124-patangiyu, GU-C124-khiskoli, GU-C124-marghi, GU-C124-gal, GU-C124-kapal

\chapter{Butterfly, Squirrel, Hen, Cheek, Forehead}
\label{ch:gu-butterfly-squirrel-hen-cheek-forehead}
\hlchaptermodality{124}

\begin{chapteropening}
\textbf{By the end of this chapter:} \emph{I can say a butterfly, a squirrel, a hen, a cheek and a forehead.}

Complete the last lesson of chapter 124: i can say a butterfly, a squirrel, a hen, a cheek and a forehead.
\end{chapteropening}

\section[\texorpdfstring{\gu{પતંગિયું} (patangiyũ)}{પતંગિયું (patangiyũ)}]{\gu{પતંગિયું} (patangiyũ) — a butterfly}

\label{lesson:GU-C124-patangiyu}



\begin{quote}
Before the new one: say the Gujarati for a fly, then the Gujarati for a bee.
\end{quote}

\subsection*{You'll want to know: \gu{પતંગિયું}}
\textbf{\gu{પતંગિયું}} (\emph{patangiyũ}) — \textquotedblleft{}a butterfly\textquotedblright{}.

\begin{grammarlens}[title={What you've built}]
One more word for this chapter.
\end{grammarlens}

\subsection*{Your turn}
\begin{itemize}
\raggedright
  \item \emph{Say it:} \emph{patangiyũ}
  \item \emph{Say it:} \emph{patangiyũ}, once more
  \item \emph{Say it:} say \emph{patangiyũ}
  \item \emph{Recall:} say the Gujarati for a fly, then the Gujarati for a bee, then say \emph{patangiyũ} again
\end{itemize}

\begin{tcolorbox}[breakable,colback=teal!4,colframe=teal!35!black,title={Before you move on}]
What does \gu{પતંગિયું} mean? (\textquotedblleft{}A butterfly\textquotedblright{}.) Say it once more.
\end{tcolorbox}
\section[\texorpdfstring{\gu{ખિસકોલી} (khiskolī)}{ખિસકોલી (khiskolī)}]{\gu{ખિસકોલી} (khiskolī) — a squirrel}

\label{lesson:GU-C124-khiskoli}



\begin{quote}
Before the new one: say the Gujarati for a bee, then the Gujarati for a butterfly.
\end{quote}

\subsection*{You'll want to know: \gu{ખિસકોલી}}
\textbf{\gu{ખિસકોલી}} (\emph{khiskolī}) — \textquotedblleft{}a squirrel\textquotedblright{}.

\begin{grammarlens}[title={What you've built}]
One more word for this chapter.
\end{grammarlens}

\subsection*{Your turn}
\begin{itemize}
\raggedright
  \item \emph{Say it:} \emph{khiskolī}
  \item \emph{Say it:} \emph{khiskolī}, once more
  \item \emph{Say it:} say \emph{khiskolī}
  \item \emph{Recall:} say the Gujarati for a bee, then the Gujarati for a butterfly, then say \emph{khiskolī} again
\end{itemize}

\begin{tcolorbox}[breakable,colback=teal!4,colframe=teal!35!black,title={Before you move on}]
What does \gu{ખિસકોલી} mean? (\textquotedblleft{}A squirrel\textquotedblright{}.) Say it once more.
\end{tcolorbox}
\section[\texorpdfstring{\gu{મરઘી} (marghī)}{મરઘી (marghī)}]{\gu{મરઘી} (marghī) — a hen}

\label{lesson:GU-C124-marghi}



\begin{quote}
Before the new one: say the Gujarati for a butterfly, then the Gujarati for a squirrel.
\end{quote}

\subsection*{You'll want to know: \gu{મરઘી}}
\textbf{\gu{મરઘી}} (\emph{marghī}) — \textquotedblleft{}a hen\textquotedblright{}.

\begin{grammarlens}[title={What you've built}]
One more word for this chapter.
\end{grammarlens}

\subsection*{Your turn}
\begin{itemize}
\raggedright
  \item \emph{Say it:} \emph{marghī}
  \item \emph{Say it:} \emph{marghī}, once more
  \item \emph{Say it:} say \emph{marghī}
  \item \emph{Recall:} say the Gujarati for a butterfly, then the Gujarati for a squirrel, then say \emph{marghī} again
\end{itemize}

\begin{tcolorbox}[breakable,colback=teal!4,colframe=teal!35!black,title={Before you move on}]
What does \gu{મરઘી} mean? (\textquotedblleft{}A hen\textquotedblright{}.) Say it once more.
\end{tcolorbox}
\section[\texorpdfstring{\gu{ગાલ} (gāl)}{ગાલ (gāl)}]{\gu{ગાલ} (gāl) — a cheek}

\label{lesson:GU-C124-gal}



\begin{quote}
Before the new one: say the Gujarati for a squirrel, then the Gujarati for a hen.
\end{quote}

\subsection*{You'll want to know: \gu{ગાલ}}
\textbf{\gu{ગાલ}} (\emph{gāl}) — \textquotedblleft{}a cheek\textquotedblright{}.

\begin{grammarlens}[title={What you've built}]
One more word for this chapter.
\end{grammarlens}

\subsection*{Your turn}
\begin{itemize}
\raggedright
  \item \emph{Say it:} \emph{gāl}
  \item \emph{Say it:} \emph{gāl}, once more
  \item \emph{Say it:} say \emph{gāl}
  \item \emph{Recall:} say the Gujarati for a squirrel, then the Gujarati for a hen, then say \emph{gāl} again
\end{itemize}

\begin{tcolorbox}[breakable,colback=teal!4,colframe=teal!35!black,title={Before you move on}]
What does \gu{ગાલ} mean? (\textquotedblleft{}A cheek\textquotedblright{}.) Say it once more.
\end{tcolorbox}
\section[\texorpdfstring{\gu{કપાળ} (kapāḷ)}{કપાળ (kapāḷ)}]{\gu{કપાળ} (kapāḷ) — a forehead}

\label{lesson:GU-C124-kapal}



\begin{quote}
Before the new one: say the Gujarati for a hen, then the Gujarati for a cheek.
\end{quote}

\subsection*{You'll want to know: \gu{કપાળ}}
\textbf{\gu{કપાળ}} (\emph{kapāḷ}) — \textquotedblleft{}a forehead\textquotedblright{}.

\begin{grammarlens}[title={What you've built}]
One more word for this chapter.
\end{grammarlens}

\subsection*{Your turn}
\begin{itemize}
\raggedright
  \item \emph{Say it:} \emph{kapāḷ}
  \item \emph{Say it:} \emph{kapāḷ}, once more
  \item \emph{Say it:} say \emph{kapāḷ}
  \item \emph{Recall:} say the Gujarati for a hen, then the Gujarati for a cheek, then say \emph{kapāḷ} again
\end{itemize}

\begin{tcolorbox}[breakable,colback=teal!4,colframe=teal!35!black,title={Before you move on}]
What does \gu{કપાળ} mean? (\textquotedblleft{}A forehead\textquotedblright{}.) Say it once more.
\end{tcolorbox}
